\chapter{Risque de modèle et validation}\label{ch:rc:model-risk-and-validation}

Le 30 janvier 2012, le groupe de revue des modèles de JPMorgan a autorisé un nouveau modèle de \omterm{def:rc:market-risk-measures:var}{valeur en risque} (VaR) pour le
portefeuille de crédit synthétique du Chief Investment Office de la banque. L'ancien modèle avait été jugé
trop conservateur~; le nouveau modèle a réduit de moitié la VaR du service, de 132 à 66 millions de dollars, et la
limite de VaR à l'échelle de la banque, dépassée pendant plusieurs jours, n'était plus franchie. Quelques mois plus tard, après que les pertes sur le portefeuille eurent atteint des milliards de dollars, le
groupe d'enquête interne de la banque a relevé une erreur opérationnelle dans les tableurs du modèle~: après avoir soustrait
une ancienne intensité de défaut ou une ancienne corrélation d'une nouvelle, le tableur divisait par leur somme au lieu de
leur moyenne, ce qui atténuait probablement la volatilité d'un facteur deux et abaissait la VaR. Le modèle avait
été approuvé~; personne n'avait vérifié qu'il calculait bien ce que disait sa documentation. Ce chapitre traite
de la discipline qui s'en charge~: identifier les modèles, les hiérarchiser selon le préjudice qu'ils peuvent
causer, les valider par rapport à des références indépendantes et aux résultats observés, et assurer la gouvernance des résultats.

\section{Ce qu'est le risque de modèle et par où il s'introduit}

\begin{definition}[Risque de modèle]\label{def:rc:model-risk-and-validation:risk}
Le \emph{risque de modèle}\index{risque de modèle} (model risk) est la possibilité de conséquences défavorables découlant de décisions
fondées sur les résultats d'un modèle~: perte financière, présentation erronée des comptes ou mauvaises décisions de risque. Il s'introduit par
une théorie ou des hypothèses erronées, des erreurs d'implémentation et de données, et l'usage d'un modèle solide hors
des conditions pour lesquelles il a été conçu.
\end{definition}

\begin{definition}[Gestion du risque de modèle]\label{def:rc:model-risk-and-validation:mrm}
La \emph{gestion du risque de modèle}\index{gestion du risque de modèle} (model risk management) est le cadre qui recense les modèles
de l'établissement, évalue et limite le risque qu'ils portent, et attribue les responsabilités liées au développement, à l'usage,
à la validation et au contrôle tout au long de leur cycle de vie.
\end{definition}

Les trois sources d'erreur appellent des contrôles différents. Une théorie erronée est détectée par une revue conceptuelle et par
la comparaison avec d'autres modèles~; les erreurs d'implémentation par une réimplémentation indépendante et par
des tests sur des cas à solution connue~; l'usage abusif par le suivi des résultats du modèle et de ses entrées.

\section{Attentes des superviseurs}

\begin{definition}[Remise en cause effective]\label{def:rc:model-risk-and-validation:challenge}
La \emph{remise en cause effective}\index{remise en cause effective} (effective challenge) est l'analyse critique d'un modèle par des personnes objectives
possédant l'expertise nécessaire pour en trouver les faiblesses, l'indépendance pour les signaler et l'autorité
nécessaire pour les faire corriger.
\end{definition}

\begin{definition}[Trois lignes de défense]\label{def:rc:model-risk-and-validation:lines}
Les \emph{trois lignes de défense}\index{trois lignes de défense} (three lines of defence) attribuent le \omterm{def:rc:model-risk-and-validation:risk}{risque de modèle} aux
propriétaires et aux utilisateurs du modèle (première ligne~: développement et usage), à une fonction de risque indépendante (deuxième ligne~:
validation et contrôle) et à l'audit interne (troisième ligne~: assurance que le dispositif fonctionne).
\end{definition}

\begin{dated}{2026-09}{Orientations sur le risque de modèle}\label{dat:rc:model-risk-and-validation:rules}
Aux États-Unis, la Fed, l'OCC et la FDIC ont publié des orientations révisées sur la \omterm{def:rc:model-risk-and-validation:mrm}{gestion du risque
de modèle} le 17 avril 2026 (SR 26-2), en remplacement des orientations d'avril 2011 (SR 11-7). Elles concernent surtout
les organisations bancaires dont les actifs dépassent 30~milliards de dollars~; elles définissent un modèle comme une
méthode quantitative complexe appliquant des théories statistiques, économiques ou financières pour transformer des
données d'entrée en estimations, à l'exclusion des calculs arithmétiques simples, comme les calculs sur tableur, et des
procédures déterministes fondées sur des règles, et elle laisse hors de son champ l'IA générative et agentique. Au Royaume-Uni, la PRA
dans sa déclaration de supervision SS1/23 (publiée en mai 2023) s'applique depuis le 17 mai 2024 aux banques disposant de modèles internes
pour le capital réglementaire, selon cinq principes~: identification et classification, gouvernance,
développement et usage, validation indépendante et mesures d'atténuation du risque. Dans la zone euro, la BCE a révisé son
guide des modèles internes le 28 juillet 2025, en ajoutant des attentes sur les techniques d'apprentissage automatique.
\end{dated}

\section{Le rapport de validation}

\begin{definition}[Validation des modèles]\label{def:rc:model-risk-and-validation:validation}
La \emph{validation des modèles}\index{validation des modèles} (model validation) est l'ensemble des activités qui vérifient qu'un
modèle fonctionne comme prévu~: solidité conceptuelle, vérification de l'implémentation, comparaison aux \omterm{def:rc:model-risk-and-validation:bench}{modèles de référence} et \omterm{def:rc:model-risk-and-validation:bench}{analyse
des résultats}, avec les constats, leur gravité et les conditions d'usage.
\end{definition}

\begin{definition}[Modèle de référence, analyse des résultats]\label{def:rc:model-risk-and-validation:bench}
Un \emph{modèle de référence}\index{modèle de référence} (benchmark) est une implémentation indépendante, du même
modèle ou d'une alternative, à laquelle on compare les sorties d'un modèle. \emph{Analyse des
résultats}\index{analyse des résultats} (outcomes analysis) compare les sorties d'un modèle aux résultats réalisés qu'ils
prédisaient~: backtest de la VaR et de la marge, attribution du P\&L, taux de défaut comparés aux PD.
\end{definition}

\begin{method}[Comparaison sur une grille]\label{met:rc:model-risk-and-validation:grid}
Choisir les plages de paramètres que le modèle rencontrera à l'usage. Évaluer le modèle et le \omterm{def:rc:model-risk-and-validation:bench}{modèle de référence} sur
chaque combinaison. Fixer les tolérances en termes opérationnels (une erreur absolue en points de base de notionnel,
et une erreur relative). Consigner chaque point, résumer les échecs par paramètre, et examiner la
tendance qui s'en dégage, et non seulement le pire cas.
\end{method}

\omcode{code/firm/modelval/firm_modelval.py}{57}{74}{Banc d'essai de comparaison~: chaque point de la grille est consigné, avec des tolérances absolues et relatives, et les échecs sont résumés.}\label{lst:rc:model-risk-and-validation:bench}

\begin{example}[Validation de l'arbre du chapitre~7]\label{ex:rc:model-risk-and-validation:tree}
L'arbre trinomial de Hull--White du chapitre~7 évalue des swaptions européennes~; la forme fermée de Jamshidian du
même modèle sert de référence. Sur 216 points (retour à la moyenne 1~\%, 5~\%, 20~\%~; volatilité 60 et 120
points de base~; échéances de 1, 5 et 10 ans~; maturités de 5 et 10 ans~; strikes à la monnaie et $\pm1~\%$~; pas trimestriels et
mensuels), avec une tolérance d'un demi-point de base de notionnel ou de 1~\%~: avec des pas trimestriels, 54 points sur 108
échouent et la pire erreur est de 15.8 points de base de notionnel~; avec des pas mensuels, 23 échouent et la pire
erreur est de 3.2. Les échecs se concentrent aux échéances courtes, 14 des 23 à un an
(la \cref{fig:rc:model-risk-and-validation:validation})~: un payeur à la monnaie, d'échéance un an sur un swap de cinq ans
(retour à la moyenne 1~\%, volatilité 60 points de base) vaut 104.81 points de base en forme fermée, 103.07 dans
l'arbre mensuel et 98.55 dans l'arbre à pas trimestriels. Le constat~: l'arbre a besoin de pas plus fins avant une date
d'exercice proche, et une condition d'usage tant qu'il n'en dispose pas.
\end{example}

\begin{omfigure}
\begin{tikzpicture}
\begin{axis}[width=11.5cm,height=5.2cm,ybar,bar width=14pt,
  symbolic x coords={1y,5y,10y},xtick=data,xticklabels={1A,5A,10A},enlarge x limits=0.3,
  xlabel={échéance de la swaption},ylabel={erreur maximale (pb de notionnel)},
  tick label style={font=\small},label style={font=\small},
  ymin=0,ymax=18,grid=major,grid style={gray!25},
  nodes near coords,every node near coord/.append style={font=\footnotesize,/pgf/number format/fixed,/pgf/number format/precision=1},
  legend columns=2,legend style={font=\footnotesize,at={(0.5,-0.32)},anchor=north,draw=none,fill=none,/tikz/every even column/.append style={column sep=8pt}},
  table/col sep=comma]
\addplot[fill=omThm!60,draw=omThm] table[x=expiry,y=err_q]{figdata/rates-credit-risk/26-model-risk-and-validation/validation.csv};
\addplot[fill=omDef!60,draw=omDef] table[x=expiry,y=err_m]{figdata/rates-credit-risk/26-model-risk-and-validation/validation.csv};
\legend{pas trimestriels,pas mensuels}
\end{axis}
\end{tikzpicture}

\omcaption{Erreur absolue maximale de l'arbre trinomial par rapport à la forme fermée, par échéance et taille du
pas, sur la grille de validation. Affiner le pas divise l'erreur par un facteur d'environ trois à cinq, surtout
aux échéances courtes, où la date d'exercice ne tombe qu'après quelques pas. Source~:
tutoriel du chapitre.}\label{fig:rc:model-risk-and-validation:validation}
\end{omfigure}

\section{Inventaire, hiérarchisation et limites}

\begin{definition}[Inventaire des modèles, hiérarchisation des modèles]\label{def:rc:model-risk-and-validation:inventory}
Un \emph{inventaire des modèles}\index{inventaire des modèles} (model inventory) recense chaque modèle en usage avec son propriétaire, sa finalité,
ses utilisateurs, son statut de validation et ses limites connues. \emph{hiérarchisation des modèles}\index{hiérarchisation des modèles} (tiering) classe
les modèles selon le dommage qu'ils peuvent causer (importance de l'usage, complexité, incertitude
des entrées et des hypothèses) et fixe en conséquence la profondeur et la fréquence de la validation.
\end{definition}

\begin{example}[L'inventaire de l'entreprise]\label{ex:rc:model-risk-and-validation:inventory}
On attribue à huit modèles de ce livre une note de 1 à 3 pour l'importance de l'usage (comptée deux fois), la complexité et
l'incertitude, en fixant le niveau~1 à partir d'un score de 10 ou plus et le niveau~2 à partir de 7 ou plus (une
politique illustrative), on obtient cinq modèles de niveau~1 (le cube SABR, le \omterm{def:rc:mortgage-modelling:model}{modèle de remboursement anticipé}, le moteur d'exposition, le modèle de VaR
et le modèle de dépôts), deux de niveau~2 (la courbe et le pricer de \omterm{def:rc:bermudans-and-callables:bermudan}{swaptions bermudéennes}) et un de niveau~3 (le
modèle de prévision des appels de marge). Les modèles de niveau~1 sont revalidés chaque année, ceux de niveau~2 tous les deux ans, ceux de niveau~3 tous les trois ans.
\end{example}

\begin{definition}[Ajustement post-modèle]\label{def:rc:model-risk-and-validation:pma}
Un \emph{ajustement post-modèle}\index{ajustement post-modèle} (post-model adjustment) est une modification documentée
des sorties d'un modèle, faite en dehors de celui-ci pour corriger une faiblesse connue jusqu'à ce que le modèle soit corrigé~; il requiert un propriétaire, une
justification, une taille et une échéance.
\end{definition}

\section{Comparaison à un modèle de référence et analyse des résultats}

La comparaison à un \omterm{def:rc:model-risk-and-validation:bench}{modèle de référence} révèle les erreurs de mise en œuvre que la documentation masque. L'erreur de 2012 en est un cas~: la
documentation indiquait des variations relatives, le tableur en calculait la moitié, et aucune mise en œuvre indépendante
ne les a comparées.

\begin{example}[La moitié de la volatilité]\label{ex:rc:model-risk-and-validation:half}
Un portefeuille de crédit synthétique illustratif vend 1~milliard d'euros de protection sur la tranche equity 0--3~\% dans le
modèle à grand panier et achète 12.0 fois plus de protection sur indice comme couverture (intensité de défaut 1~\%, corrélation 20~\%).
Sur une année de variations quotidiennes illustratives de l'intensité de défaut et de la corrélation, sa VaR historique à 99~\% est
de 14.06~millions d'euros avec des variations relatives divisées par la moyenne des anciennes et des nouvelles valeurs, et de 7.01~millions
d'euros lorsqu'elles sont divisées par la somme~: 49.9~\% de la valeur visée
(\cref{fig:rc:model-risk-and-validation:losses}). Un \omterm{def:rc:model-risk-and-validation:bench}{modèle de référence} qui recalculerait un scénario à la main
l'aurait relevé.
\end{example}

\begin{omfigure}
\begin{tikzpicture}
\begin{axis}[width=11.5cm,height=5.2cm,
  xlabel={rang de la perte du scénario},ylabel={perte (millions d'euros)},
  tick label style={font=\small},label style={font=\small},
  xmin=1,xmax=40,ymin=0,ymax=16,grid=major,grid style={gray!25},
  legend columns=2,legend style={font=\footnotesize,at={(0.5,-0.25)},anchor=north,draw=none,fill=none,/tikz/every even column/.append style={column sep=8pt}},
  table/col sep=comma]
\addplot[omDef,very thick,mark=*,mark size=1pt] table[x=rank,y=intended]{figdata/rates-credit-risk/26-model-risk-and-validation/losses.csv};
\addplot[omThm,very thick,mark=square*,mark size=1pt] table[x=rank,y=error]{figdata/rates-credit-risk/26-model-risk-and-validation/losses.csv};
\draw[omExo,dashed] (axis cs:3,0) -- (axis cs:3,16);
\legend{variations divisées par la moyenne,variations divisées par la somme}
\end{axis}
\end{tikzpicture}

\omcaption{Les quarante plus fortes pertes de scénario du portefeuille de crédit synthétique illustratif, avec des variations relatives
calculées comme prévu et comme dans le tableur de 2012. La ligne en pointillés marque le scénario qui fixe
la VaR à 99~\%~: chaque perte et la VaR sont à peu près divisées par deux. Source~: historique illustratif~; le
tutoriel du chapitre.}\label{fig:rc:model-risk-and-validation:losses}
\end{omfigure}

\begin{dated}{2026-09}{Le modèle de 2012}\label{dat:rc:model-risk-and-validation:cio}
Selon le rapport du groupe de travail sur la gestion de JPMorgan (janvier 2013), le nouveau modèle de VaR pour le Chief
Investment Office, portefeuille de crédit synthétique, a été autorisé le 30 janvier 2012 et a mis fin aux dépassements
de la limite de VaR à l'échelle de l'entreprise~; parmi les erreurs relevées ensuite figurait un tableur qui divisait la
variation d'une intensité de défaut ou d'une corrélation par la somme des anciennes et des nouvelles valeurs au lieu de leur moyenne,
réduisant vraisemblablement la volatilité d'un facteur deux. Le rapport a aussi constaté que les tests de prix reposaient sur
des tableurs qui n'avaient pas été vérifiés.
\end{dated}

L'\omterm{def:rc:model-risk-and-validation:bench}{analyse des résultats} détecte ce que la comparaison ne peut pas détecter~: un modèle mis en œuvre conformément à sa spécification mais
faux sur le monde réel. Ses outils sont les backtests des chapitres~21 et~25, le test d'attribution du
chapitre~23 et le suivi des entrées par rapport aux plages validées.

\section{Tutoriel~: un rapport de validation}\label{tut:rc:model-risk-and-validation:tutorial}

\begin{tutorial}
\textbf{Objectif.} Valider l'arbre du chapitre~7 par rapport à la forme fermée sur une grille, consigner les échecs,
hiérarchiser les modèles de l'entreprise et rejouer l'erreur de moyenne. \textbf{État final~:} les chiffres des
\cref{ex:rc:model-risk-and-validation:tree,ex:rc:model-risk-and-validation:inventory,ex:rc:model-risk-and-validation:half}
et les graphiques, rédigés selon la structure ci-dessous.
\begin{tutsteps}
\item \textbf{Périmètre}~: le modèle, ses usages, les plages de paramètres en vigueur (\texttt{GRID}).
\item \textbf{Comparaison}~: \texttt{validation()} exécute \texttt{benchmark} avec \texttt{tree\_price} face à \texttt{closed\_price}.
\item \textbf{Constats}~: \texttt{errors\_by\_expiry()}~; gravité et conditions d'utilisation.
\item \textbf{Inventaire}~: \texttt{inventory()}, \texttt{tiers()}~; \texttt{averaging\_error()}~;
\texttt{fig\_rc\_modelval.py} écrit les graphiques.
\end{tutsteps}
\textbf{Pour aller plus loin.} Ajouter une comparaison Monte-Carlo des mêmes swaptions~; ajouter une grille de pas plus fine
au voisinage de la date d'exercice et relancer les échecs sur un an.
\end{tutorial}

\section{Construction~: outils de validation des modèles}\label{bld:rc:model-risk-and-validation:modelval}

\begin{build}
\bfield{Rôle} L'inventaire et l'outillage de validation de l'entreprise, utilisés pour chaque composant de cette série
avant son entrée dans le moteur de risque.
\bfield{Interface} \texttt{ModelRecord} (avec \texttt{tier}, \texttt{revalidate\_years})~;
\texttt{inventory\_summary}~; \texttt{benchmark(candidate, reference, grid, abs\_tol, rel\_tol)},
\texttt{BenchmarkResult}, \texttt{summary}~; \texttt{binomial\_tail}~; \texttt{relative\_change}.
\bfield{Règles} Chaque point de grille est conservé, et non pas seulement les échecs~; les tolérances sont exprimées en unités métier~; les scores de hiérarchisation et
les intervalles sont un paramètre de politique.
\bfield{Tests d'acceptation} \texttt{code/firm/modelval/tests/}~: bornes de hiérarchisation~; le harnais repère le
seul échec inséré~; la queue binomiale comparée au tableau de Bâle~; les deux variations relatives diffèrent d'un facteur
deux.
\bfield{Pour aller plus loin} Rapports de validation générés à partir des enregistrements~; suivi des plages d'entrée en production~;
un outil de suivi des constats avec propriétaires et échéances.
\end{build}

\begin{omsources}
\item JPMorgan Chase \& Co., \emph{Report of JPMorgan Chase \& Co. Management Task Force Regarding 2012
CIO Losses}, 16 janvier 2013.
\item Board of Governors of the Federal Reserve System, OCC and FDIC, \emph{Revised Guidance on Model
Risk Management}, SR 26-2, 17 avril 2026.
\item Prudential Regulation Authority, SS1/23, \emph{Model risk management principles for banks}, 2023.
\end{omsources}

\section{Exercices}

\begin{exercise}[$\star$]\label{exo:rc:model-risk-and-validation:1}
Un taux passe de 1.00~\% à 1.10~\%. Donner la variation relative telle que voulue et telle que calculée dans le
tableur de 2012.
\end{exercise}

\begin{exercise}[$\star$]\label{exo:rc:model-risk-and-validation:2}
Donner le niveau de hiérarchisation d'un modèle noté 3 sur la matérialité, 1 sur la complexité et 1 sur
l'incertitude selon la règle du chapitre.
\end{exercise}

\begin{exercise}[$\star$]\label{exo:rc:model-risk-and-validation:3}
Un modèle de VaR présente 8 dépassements en 250 jours à 99~\%. Quelle est la probabilité d'un tel résultat pour un modèle correct~?
\end{exercise}

\begin{exercise}[$\star\star$]\label{exo:rc:model-risk-and-validation:4}
Pourquoi une erreur de tableur qui divise la volatilité par deux ne serait-elle pas détectée par un backtest de VaR en quelques semaines~?
\end{exercise}

\begin{exercise}[$\star\star$]\label{exo:rc:model-risk-and-validation:5}
Pourquoi une forme fermée du même modèle est-elle un \omterm{def:rc:model-risk-and-validation:bench}{modèle de référence} (benchmark) plus faible qu'un modèle alternatif, et pourquoi reste-t-elle
néanmoins utile~?
\end{exercise}

\begin{exercise}[$\star\star$]\label{exo:rc:model-risk-and-validation:6}
Selon la définition de SR 26-2, le tableur de 2012 est-il un modèle~? Quelles conséquences en découlent pour son contrôle~?
\end{exercise}

\begin{exercise}[$\star\star\star$]\label{exo:rc:model-risk-and-validation:7}
\emph{Programmation.} Relancer la validation de l'arbre avec des pas d'un vingt-quatrième d'année pour les
échéances d'un an uniquement, et signaler les échecs.
\end{exercise}

\begin{exercise}[$\star\star\star$]\label{exo:rc:model-risk-and-validation:8}
\emph{Trouver la faille.} «~Le nouveau modèle a été approuvé par le groupe de revue, la VaR plus basse est donc le bon
chiffre.~»
\end{exercise}

\section{Problème~: L'erreur de moyenne}

\begin{problem}[{Problème du week-end --- un facteur deux}]\label{pb:rc:model-risk-and-validation:1}
Le portefeuille de crédit synthétique illustratif de l'\cref{ex:rc:model-risk-and-validation:half} est mesuré par un
modèle de VaR implanté dans un tableur qui divise les variations relatives par la somme des anciennes et des nouvelles valeurs.

\medskip\noindent\textbf{Partie I --- L'erreur.}
\begin{enumerate}
\item Montrer que l'erreur divise par deux chaque variation relative.
\item Donner la VaR avec et sans l'erreur, ainsi que le rapport.
\item Pourquoi le rapport n'est-il pas exactement égal à un demi~?
\item Donner le ratio de couverture du portefeuille et expliquer pourquoi il est si élevé.
\item Quel comportement à la limite l'erreur produirait-elle~?
\end{enumerate}

\medskip\noindent\textbf{Partie II --- Détection.}
\begin{enumerate}[resume]
\item Quelle étape de validation l'aurait détectée, et en combien de temps~?
\item Qu'aurait révélé l'\omterm{def:rc:model-risk-and-validation:bench}{analyse des résultats}, et à quel moment~?
\item Quel rôle a joué le moment du changement de modèle en 2012~?
\item Pourquoi la documentation du modèle n'a-t-elle pas révélé l'erreur~?
\item Quel contrôle sur les tableurs faudrait-il exiger~?
\end{enumerate}

\medskip\noindent\textbf{Partie III --- Gouvernance.}
\begin{enumerate}[resume]
\item Qui devrait approuver un modèle qui abaisse la VaR d'un desk pendant un dépassement de limite~?
\item Quelles conditions d'utilisation faudrait-il y attacher~?
\item Comment faut-il traiter les limites lorsqu'un modèle change~?
\item Quel niveau de hiérarchisation faudrait-il attribuer à ce modèle, et pourquoi~?
\item Qu'exige ici de la part du validateur la \omterm{def:rc:model-risk-and-validation:challenge}{remise en cause effective} (effective challenge)~?
\end{enumerate}

\medskip\noindent\textbf{Partie IV --- Jugement.}
\begin{enumerate}[resume]
\item La perte de 2012 a-t-elle été causée par le modèle~?
\item Que signifie, pour des cas de ce type, l'exclusion par SR 26-2 des simples calculs de tableur~?
\item Comment faudrait-il concevoir une validation qui détecte à moindre coût les erreurs d'implémentation~?
\item Énoncer le \emph{résultat nommé}~: la réduction de VaR produite en divisant une variation de taux par la somme de
deux taux au lieu de leur moyenne.
\item En une phrase~: qu'est-ce que le \omterm{def:rc:model-risk-and-validation:risk}{risque de modèle}~?
\end{enumerate}
\end{problem}

\section{Questions d'entretien}

\begin{interviewq}[$\star$ \iqroles{risque, banque}]\label{iq:rc:model-risk-and-validation:1}
Que contient un rapport de validation de modèle~?
\end{interviewq}

\begin{interviewq}[$\star\star$ \iqroles{chercheur, risque}]\label{iq:rc:model-risk-and-validation:2}
Comment valideriez-vous un modèle de valorisation pour les \omterm{def:rc:bermudans-and-callables:bermudan}{swaptions bermudéennes}~?
\end{interviewq}

\begin{interviewq}[$\star\star$ \iqroles{développeur}]\label{iq:rc:model-risk-and-validation:3}
Comment testez-vous un moteur de valorisation numérique pour que les erreurs d'implémentation soient détectées tôt~?
\end{interviewq}

\begin{interviewq}[$\star\star$ \iqroles{risque}]\label{iq:rc:model-risk-and-validation:4}
Qu'est-ce que la \omterm{def:rc:model-risk-and-validation:challenge}{remise en cause effective} (effective challenge), et comment savez-vous qu'elle a lieu~?
\end{interviewq}

\begin{interviewq}[$\star\star\star$ \iqroles{banque, risque}]\label{iq:rc:model-risk-and-validation:5}
Qu'est-ce qui a mal tourné avec le modèle de VaR du CIO en 2012, et quels contrôles l'auraient arrêté~?
\end{interviewq}

\begin{interviewq}[$\star\star\star$ \iqroles{chercheur}]\label{iq:rc:model-risk-and-validation:6}
Comment quantifieriez-vous le \omterm{def:rc:model-risk-and-validation:risk}{risque de modèle} dans un prix, en vue d'une provision~?
\end{interviewq}
